% Source: https://artofproblemsolving.com/community/c3h1710124_junior_regional_olympiad__fbh_2017_grade_6_problem_2
% Stable identity: aops-topic-1710124-q01
Let the amount remaining in each cistern be $t$. Then the first cistern held $\dfrac{3}{4}$ of its milk equal to $t$, so it held $\dfrac{4t}{3}$; similarly the second held $\dfrac{5t}{4}$ (since $\dfrac{4}{5}$ of it is $t$) and the third $\dfrac{7t}{4}$ (since $\dfrac{4}{7}$ of it is $t$). Then
\[
\dfrac{4t}{3}+\dfrac{5t}{4}+\dfrac{7t}{4}=\dfrac{16t+15t+21t}{12}=\dfrac{13t}{3}=780,
\]
so $t=180$. The cisterns hold $\boxed{240,\ 225\text{ and }315}$ litres. Check: $180=\dfrac{3}{4}\cdot240=\dfrac{4}{5}\cdot225=\dfrac{4}{7}\cdot315$, and $240+225+315=780$.
